\chapter{Base change, reduced fibers, and permanence}\label{ch:base-change-reduced-fibers-and-permanence}
\chaptermark{Base change and permanence}

We want to understand the behaviour of models under base change with respect to an extension of the base field.
Usually, $L|K$ will denote a finite field extension.
Since $v_K$ is Henselian, it has a unique extension to $v_L$ to $L$, and then $(L,v_L)$ also satisfies Assumption \ref{ass:v_K}.
Given a model $\X$ of $X$, we ask ourselves whether the base change $\X\times_{\OO_K}\OO_L$ is again a model of $X_L \coloneqq X\times_K L$.
Well, almost: All properties required of a model from Definition \ref{def:model} hold except possibly normality.
We obtain a model of $X_L$, called the {\em normalised base change} of $\X$ relative to $L|K$, by normalizing $\X\times_{\OO_K}\OO_L$.\footnote{
    Here we use the fact that Assumption \ref{ass:v_K} implies that $\OO_K$ is an \emph{excellent} ring.
    It follows that $\X\times_{\OO_K}\OO_L$ is an excellent scheme, and therefore its normalization is a finite morphism.
    See \cite[Chapter 8, in particular Theorem 8.2.39 (d)]{Liu}.
}
We say that a model $\X$ is \textit{permanent} if $\X \times_{\OK}\OL$ is itself a model of $X$:

\begin{definition} \label{def:permanent-model}
    A model $\X$ of $X$ is called \textit{permanent}, if $\X \times_{\OK} \OL$ is normal for all finite field extensions $L|K$.
\end{definition}

We have the following theorem about permanence:

\begin{theorem} \label{thm:permanence}
    \begin{enumerate}[(i)]
        \item For a model $\X$ of $X$,  the following three statements are equivalent:
            \begin{enumerate}[(a)]
                \item $\X$ is permanent,
                \item the special fiber $\X_s$ of $\X$ is reduced,
                \item all $v\in V(\X)$ are \emph{weakly unramified} over $v_K$, i.e.\ $v$ and $v_K$ have the same value group.
            \end{enumerate}
        \item Given a model $\X$, there exists a finite extension $L|K$ such that the normalised base change of $\X$ relative to $L|K$ is permanent.
    \end{enumerate}
\end{theorem}

In Section \ref{sec:models-with-reduced-special-fibre}, we will give a proof of part (i), see Corollary \ref{cor:permanence-part-i}.
In Section \ref{sec:models-under-base-field-extensions}, we will prove part (ii) in the special case where all mulitplicities of $\X_s$
are prime to the characteristic of the residue field $k$ (Lemma \ref{lem:abhyanka}).
In the next chapter, we see proofs of part (ii) in other important special cases.

\section{Models with Reduced Special Fibre}\label{sec:models-with-reduced-special-fibre}

As in the last two sections, we assume $X$ to be a smooth, projective and absolutely irreducible curve over the local
field $K$.


\begin{remark}
    Since flatness and properness are preserved under base change, a model $\X$ of $X$ is permanent if and only if
    $\X \times_{\OK} \OL$ is a model of $X_L = X \times_K L$.
\end{remark}

\begin{definition}
    A \textit{pre-model} of $X$ is a flat and proper $\OK$-scheme $\X$ with $\X \times_{\OK} K = X$.
    That is, $\X$ satisfies all the conditions of a model except of being normal.
\end{definition}

\begin{lemma}[{\cite[Lemma 1.3]{ArzdorfWewers}}]\label{lem:conditions-for-normality}
    Let $\X$ be a pre-model of $X$.
    Then $\X$ is normal (i.e.\ a model) if and only if the following two conditions hold.
    \begin{enumerate}[(i)]
        \item The special fibre $\X_s$ has no embedded points, i.e.\ no point $x\in \X_s$ \label{lem:conditions-for-normality-i}
            satisfying $\m_x = \Ann{R}{r}$ for some $r\in R \coloneqq \OO_{\X_s,x}$ (see~\cite[Definition 7.1.6]{Liu}).
        \item The local ring $\OO_{\X,\xi}$ is a discrete valuation ring for every generic point $\xi \in \X_s$. \label{lem:conditions-for-normality-ii}
    \end{enumerate}
\end{lemma}

\begin{proof}
    According to Serre's criterion for normality~\cite[Theorem 8.2.23]{Liu},
    we need to check the following properties:
    \begin{itemize}
        \item[$(\mathrm{R}_1)$] $\X$ is regular at $x$ when $\codim{x} \leq 1$, and
        \item[$(\mathrm{S}_2)$] $\depth{\OO_{\X,x}} \geq 2$
    \end{itemize}
    for all $x\in \X$. \\
    The points on $\X$ with codimension $1$ are precisely the points on the generic fibre, which is regular by assumption,
    and the generic point on the special fibre.
    Since any discrete valuation ring is a local Noetherian domain and vice versa, $(\mathrm{R}_1)$ is
    equivalent to (\ref{lem:conditions-for-normality-ii}).
    By~\cite[Proposition 8.2.11]{Liu}, condition $(\mathrm{S}_2)$ is equivalent to say
    that $\depth{\OO_{\X_s, x}} \geq 1$ for all $x \in \X_s$.
    But this is always true for the generic point, and if $x$ is a different point it is equivalent to say that $x$ is not an embedded point.
\end{proof}

\begin{example}
    \begin{enumerate}
        \item Consider a pre-model with special fibre
            \begin{align}
                \X_s \coloneqq \Spec{k[x,y]/(x^2,xy)},
            \end{align}
            which is topologically equivalent to $\AA^1(k)$.
            The point corresponding to the maximal ideal $(x,y)$ is an embedded point.
        \item We consider Example \textrm{III}.9.8.4 from \cite{Hartshorne} of a pre-model of a smooth curve with embedded point on the special fibre.
            Let $C$ be the curve in $\PP^3(\QQ_p) = \Proj{\QQ_p[X,Y,Z,W]}$, which is given by the ideal 
            \begin{align}
                I = \left(p^2(x+1) - z^2, px(x+1) - yz, xz - py, y^2 - x^2(x+1) \right)
            \end{align}
            on the affine chart $W \neq 0$.
            $p$ is no zero divisor in $\ZZ_p[x,y,z]/I$, thus the projective scheme $\mathcal{C}/\ZZ_p$ defined by $I$
            is flat and hence a pre-model of $C$.
            On the affine chart $W \neq 0$, the special fibre $\mathcal{C}_s$ is given by the ideal 
            \begin{align}
                I_p = (z^2, yz, xz, y^2 - x^2(x+1))
            \end{align}
            and has support equal to the nodal cubic curve $y^2 = x^2(x+1)$.
            At any point where $x \neq 0$, $\mathcal{C}_s$ is reduced.
            But in the stalk at the point $(0,0,0)$, we have the element $z$ with $z^2=0$, a non-zero nilpotent element.
            Thus, $(0,0,0)$ is an embedded point and $\mathcal{C}$ no model of $C$ by Lemma \ref{lem:conditions-for-normality}.
    \end{enumerate}
\end{example}

\begin{proposition}[{\cite[Proposition 2.3]{ArzdorfWewers}}] \label{prop:reducedness-implies-permanence}
    If $\X_s$ is reduced, then $\X$ is permanent.
\end{proposition}

\begin{proof}
    Let $L|K$ be a finite extension and $l|k$ the extension of the corresponding residue fields.
    The special fibre of $\X' \coloneqq \X_{\OL} = \X \times_{\OK} \OL$ is $\X'_s := \X_s \times_k l$,
    which is reduced since $k$ is perfect by assumption.
    This means $\X'_s$ satisfies condition (ii) from Lemma~\ref{lem:conditions-for-normality}. \\
    For condition (i) let $\xi \in \X'_s$ be any generic point.
    The stalk $\OO_{\X', \xi}$ then is a Noetherian local domain of dimension $1$ and it holds
    \begin{align}
        \OO_{\X'_s, \xi} = \OO_{\X', \xi}/\pi_L\OO_{\X', \xi},
    \end{align}
    where $\pi_L \in \OL$ is a uniformising parameter.
    Since $\X'_s$ is reduced, $\pi_L\OO_{\X',\xi}$ is maximal ideal which means that
    condition (i) from Lemma~\ref{lem:conditions-for-normality} is satisfied as well.
    Thus $\X'$ is normal and hence $\X$ is permanent.
\end{proof}

\begin{proposition}[{\cite[Corollary 2.2 (ii)]{ArzdorfWewers}}]\label{prop:criterion-reducedness}
    Let $\X$ be a model of $X$, then the following statements are equivalent.
    \begin{enumerate}[(i)]
        \item $\X_s$ is reduced.
        \item $\OO_{\X_s, \xi}$ is a field for every generic point $\eta \in \X_s$.
        \item Every $v_Z \in V(\X)$ is weakly unramified.
    \end{enumerate}
\end{proposition}

\begin{proof}
    Since $\X$ is normal, there are no embedded points in $\X_s$ by Lemma~\ref{lem:conditions-for-normality}.
    Since $\X_s$ has dimension one, it is reduced if and only if $\X_s$ is reduced in codimension zero
    which in turn is the case if $\OO_{\X_s, \xi}$ is a field for every generic point $\xi \in \X_s$. \\
    For the second equivalence recall that according to Proposition~\ref{prop:multiplicities-coincide},
    for any generic point $\xi\in \X_s$ and $Z = \overline{\lbrace\eta\rbrace}$ it holds
    \begin{align}
        [\Gamma_{v_Z}: \Gamma_K] = \length \OO_{\X_s, \xi}.
    \end{align}
    This implies that $\OO_{\X_s, \xi}$ is a field if and only $v_Z$ is weakly ramified for all $v_Z \in V(\X)$.
\end{proof}

\begin{corollary} \label{cor:permanence-part-i}
    Part (i) of Theorem \ref{thm:permanence} holds.
\end{corollary}

\begin{proof}
    The equivalence (b) $\iff$ (c) follows from Proposition \ref{prop:criterion-reducedness}.
    The implication (b) $\Rightarrow$ (a) is Proposition \ref{prop:reducedness-implies-permanence}, the other direction follows from the definition.
\end{proof}


\begin{example}[Curve with reduced special fibre]
    Let $K = \QQ_3$ and consider the projective curve
    \begin{align}
        X: \, y^3 = 3x^2 + 1
    \end{align}
    over $K$.
    The equation defines a model $\X$ over $\OK=\ZZ_3$ (cf. Example \ref{ex:model-of-plane-curve}).
    The reduction $\X_s/\FF_3$ is given by $0 = y^3 - 1 = (y - 1)^3$ which means $\X_s$ is not reduced. \\
    The irreducible component is given by $Z: \, y-1$; denote the corresponding valuation by $v$.
    For $z \coloneqq y - 1$ we have
    \begin{align*}
        3x^2 &= y^3 - 1 \\
            &= (z+1)^3 - 1 \\
            &= z^3 + 3z^2 + 3z,
    \end{align*}
    where the right hand side is an Eisenstein polynomial.
    This results $v(z) = \frac{1}{3}v_{\scriptscriptstyle G}(3x^2) = \frac{1}{3} \notin \ZZ = \Gamma_K$.
    Let us extend the ground field to $L = \QQ_3[\theta]$ where $\theta^3 - 3 = 0$ and apply the coordinate transformation
    $(x,y) \mapsto (x_1, 1+ \theta y_1)$.
    We then obtain a birational curve
    \begin{align}
        X_1: \, y_1^3 + \theta^2 y_1^2 + \theta y_1 = x_1^2.
    \end{align}
    Since $y_1^3 \equiv x_1^2$ (mod $\theta$), $\X_s$ is reduced.
    It holds $v_G(\theta y_1) = v_{\scriptscriptstyle}(y - 1) = \frac{1}{3} = v_L(\theta) \in \Gamma_L$, especially $v_{\scriptscriptstyle G}(y_1) = 0$.
\end{example}


\section{Models under Base Field Extensions}\label{sec:models-under-base-field-extensions}

Let $L|K$ be a finite field extension and let $\X$ be a model of $X$.
Denote by $\X' \coloneqq (X\times_{\OK} \OL)^\sim$ the normalised base change of $\X$ to $\OL$.
Let $\xi \in \X_s$ be a generic point and let $\xi' \in \X_s'$ be a generic point lying over $\xi$.
Since $\X$ (resp. $\X'$) is normal we have that $R \coloneqq \OO_{\X, \xi}$ (resp. $ \coloneqq \OO_{\X', \xi'}$)
is a discrete valuation ring dominating $\OK$ (resp. $\OL$) by
Lemma~\ref{lem:conditions-for-normality} (\ref{lem:conditions-for-normality-i}), that is
\begin{itemize}
    \item $R \supset \OO_K$ (resp. $S \supset \OL$), and
    \item $\m_R \cap \OK = \pi_K\OK$ (resp. $\m_S \cap \OL = \pi_L\OL$).
\end{itemize}
We arrive at the following diagram of fields.


\begin{figure}[H]\label{fig:valuations-and-field-extensions}
    \centering
    \resizebox{0.5\textwidth}{!}{
        \begin{tikzcd}[ampersand replacement=\&]
            {}
            \& {(F_{X_L},\{w_{i,1},\dots,w_{i,j_i}\})} \\
            {(F_X,\{v_i\})} \arrow[ru, no head]
            \& {} \\
            {}
            \& {(L(x),v_{{\scriptscriptstyle G}, L})} \arrow[uu, no head]  \\
            {(K(x),v_{{\scriptscriptstyle G}, K})} \arrow[ru, no head] \arrow[uu, no head]
            \& {} \\
            {}
            \& {(L,v_L)} \arrow[uu, no head]  \\
            {(K,v_K)} \arrow[ru, no head] \arrow[uu, no head]
            \& {} \\
        \end{tikzcd}
    }
\end{figure}

Here the $v_i$ are the valuations on $F_X$ corresponding to the generic points of the components of $\X_s$.
The transcendental element of the field extension $K(x)|K$ can be chosen such that the Gauß valuation $v_{{\scriptscriptstyle G}, K}$
on $K(x)$ is compatible with the extension of valuations. \\
$w_{i, j}$ represent the valuations corresponding to the irreducible component of $\X_s'$ whose generic point lies over
the generic point of $Z_i \subset \X_s$ which corresponds to $v_i$.

\begin{theorem}[{\cite[Proposition 2.3]{ArzdorfWewers}}]
    Let $\X$ be a model of $X$.
    Then there is a finite extension $L|K$ such that $\X' = (\X \times_{\OK}\OL)^\sim$ has a reduced special fibre.
\end{theorem}

\begin{proof}
    By~\cite[Theorem 1.0]{Epp} there exists a field extension $L|K$ such that $w_{i,j}$ is weakly unramified over $v_L$.
    According to Proposition~\ref{prop:criterion-reducedness}, the special fibre of $\X'$ is reduced.
\end{proof}

\begin{lemma}[Abhyanka]\label{lem:abhyanka}
    Let $L|K$ be a finite extension.
    Let $\lbrace v_1, \dots, v_n \rbrace = V(X)$ denote the extended valuation on the function field $F_X|K$.
    We have that $F_{X_L} = F_X L$ and with the extended valuation denoted as $w_{1,1}, \dots, w_{n, j_n}$ respectively.
    Assume that $v_i|v_{{\scriptscriptstyle G}, K}$ are tamely ramified for all $i$.
    If $e_i \coloneqq e(v_i|v_{{\scriptscriptstyle G},K})$ divides
    $e(v_{{\scriptscriptstyle G}, L}, v_{{\scriptscriptstyle G}, K}) = e(v_L|v_K)$ for all $i$, then $w_{i,j} \in V(X_L)$
    is weakly unramified over $v_{{\scriptscriptstyle G}, L}$ for all $i,j$.
\end{lemma}

\begin{proof}
    We will use the equivalence of Proposition~\ref{prop:criterion-reducedness} without further mention,
    that is the equivalence of the special fibre being reduced and the valuations being mildly ramified.
    Without loss of generality we may assume the following assertions:
    \begin{itemize}
        \item $n = 1$ and $j_1 = 1$ since the argument can be constructed inductively,
        \item $e \coloneqq e_1 = e(v_1|v_{{\scriptscriptstyle G},K}) = e(v_{{\scriptscriptstyle G},L}|v_{{\scriptscriptstyle G},K})$
            since we assume the residue field to be perfect, and
        \item $K$ contains the $e$-th root of unity, i.e. $\pi_K = \pi_K^e$, since we can always replace $K$ by a suitable
            Kummer extension without violating the other assumptions and obtain a compatible diagram regarding the valuations.
    \end{itemize}
    Since $F_{X_L} = F_X L$ we have that $[\Gamma_{w_1}: \Gamma_{v_1}] = 1$.
    By multiplicativity of the ramification indices we conclude $[\Gamma_{w_1}: \Gamma_{v_L}] = 1$.
\end{proof}

\begin{example}[Reduction after wildly ramified field extension]
    Consider the curve
    \begin{align}
        X: \, y^2 = 8x^2 - 2
    \end{align}
    defined over the field $K \coloneqq \QQ_2$.
    This example illustrates that the right choice of the field extension is crucial for having a reduced special fibre,
    since a wild ramified field extension is not unique.
    \begin{enumerate}[(i)]
        \item First let us consider the extension $L_1 = K(\theta_1)$ where $\theta_1^2 - 2 = 0$ and the coordinate transformation
            $(x,y) \mapsto (x_1, 2y_1 - \theta_1 - 2)$.
            We obtain a birational curve
            \begin{align}
                X_1: \, y_1^2 - (\theta_1 + 4)y_1 = 2x_1^3 - (2 + \theta_1)
            \end{align}
            with $v_{{\scriptscriptstyle G}, L_1}(y_1) = \frac{1}{4}$.
            We have $[\Gamma_{F_{X_1}}: \Gamma_{L_1}] = 2 \neq 1$.
        \item Let us consider the field extension $L_2 = K(\theta_2)$ where $\theta_2^2 + 2 = 0$
            and the coordinate transformation $(x,y) \mapsto (x_2, 2\theta_2 y_2 + \theta_2)$.
            We obtain a birational curve
            \begin{align}
                X_2: \, y_2^2 + y_2 = -x_2^3
            \end{align}
            which is given by an equation of AS-form, i.e.\ its special fibre is a reduced curve.
            Thus $[\Gamma_{F_{X_2}}: \Gamma_{L_2}] = 1$.
    \end{enumerate}
\end{example}
